\documentclass[11pt]{article}

\usepackage[T1]{fontenc}
\usepackage[utf8]{inputenc}
\usepackage{lmodern}
\usepackage{amsmath,amssymb,amsfonts,amsthm}
\usepackage[margin=1in]{geometry}
\usepackage{microtype}
\usepackage{xcolor}
\usepackage[colorlinks=true,linkcolor=teal,urlcolor=teal]{hyperref}

\newcommand{\F}{\mathcal F}
\newcommand{\T}{\mathcal T_1}
\newcommand{\Tr}{\operatorname{Tr}}
\newcommand{\B}{\mathrm B}
\newcommand{\ketbra}[1]{\lvert #1\rangle\!\langle #1\rvert}

\newtheorem*{lemmaD}{Lemma D.5 (rephrased)}

\title{A Short Guide to Lemma D.5\\
\large The tangent witness for the multimode Gaussian fidelity bound}
\author{}
\date{}

\begin{document}
\maketitle

\begin{abstract}
Lemma D.5 converts the product stochastic inequality of Lemma D.4 into the
sharp multimode root-fidelity bound needed for Gaussian amplification.  Its
key step is a tangent-line inequality for the square root.  The tangent is
taken at the photon-number law of the quantum-limited amplifier, and its
slope is exactly a bounded product of decreasing weights, so Lemma D.4
controls the linear correction.  This note derives that step explicitly and
separates it from the external least-noise theorem used earlier in the
dependency chain.
\end{abstract}

\section{Purpose of the lemma}

At the origin, the desired amplified state is the original product thermal
state $\Phi_0$, while every physical amplifier of gain $\gamma>1$ must add
noise.  The quantum-limited amplifier produces a hotter product thermal state
$\tau^{(\gamma)}$.  Lemma D.5 proves that no displacement- and
phase-covariant operation, even a trace-decreasing one, can have larger root
fidelity with $\Phi_0$ than this quantum-limited output.

The proof has a useful convex-analytic form.  Root fidelity with a fixed
diagonal target is a concave function of the output probabilities.  Its
tangent plane at the quantum-limited distribution is a linear
photon-number test.  Earlier stochastic-order lemmas show that every allowed
output lies on the correct side of that test.

\section{Background}

\subsection{Root fidelity for diagonal positive operators}

For positive trace-class operators $R,S$, possibly subnormalized, the paper
uses root fidelity
\[
  \B(R,S)=\Tr\sqrt{R^{1/2}SR^{1/2}}.
\]
If $R$ and $S$ are diagonal in the same basis, with masses
$P=(P_{\mathbf k})$ and $t=(t_{\mathbf k})$, then
\begin{equation}\label{eq:classical-B}
  \B(R,S)=\B(P,t):=
  \sum_{\mathbf k}\sqrt{P_{\mathbf k}t_{\mathbf k}}.
\end{equation}
This is the classical Hellinger affinity.  The formula remains valid when
$\sum_{\mathbf k}P_{\mathbf k}<1$.

For tensor products, root fidelity is multiplicative:
\begin{equation}\label{eq:multiplicative}
  \B\!\left(\bigotimes_\ell R_\ell,
             \bigotimes_\ell S_\ell\right)
  =\prod_\ell\B(R_\ell,S_\ell).
\end{equation}

\subsection{Thermal laws and the one-mode factor}

For $0<q<1$,
\[
  \tau_q=(1-q)\sum_{k\geq0}q^k\ketbra{k}.
\]
At intensity gain $\gamma>1$, define
\begin{equation}\label{eq:qgamma}
  q^{(\gamma)}=1-\frac{1-q}{\gamma}.
\end{equation}
The quantum-limited amplifier sends $\tau_q$ to
$\tau_{q^{(\gamma)}}$.  Summing a geometric series gives
\begin{align}
  b_\gamma(q)
  &:=\B(\tau_{q^{(\gamma)}},\tau_q)\notag\\
  &=\frac{\sqrt{(1-q)(1-q^{(\gamma)})}}
          {1-\sqrt{q q^{(\gamma)}}}
   =\frac{\sqrt\gamma+\sqrt{q(\gamma-1+q)}}{\gamma+q}.
  \label{eq:bgamma}
\end{align}
The last expression follows by substituting \eqref{eq:qgamma} and
rationalizing the denominator.

For $r$ modes, set
\begin{equation}\label{eq:product-states}
  \Phi_0=\bigotimes_{\ell=1}^r\tau_{q_\ell},
  \qquad
  \tau^{(\gamma)}
  =\bigotimes_{\ell=1}^r\tau_{q_\ell^{(\gamma)}}.
\end{equation}
Their joint number distributions are
\begin{align}
  t_{\mathbf k}
  &=\prod_{\ell=1}^r(1-q_\ell)q_\ell^{k_\ell},
  \label{eq:t}\\
  t_{\mathbf k}^{(\gamma)}
  &=\prod_{\ell=1}^r
      (1-q_\ell^{(\gamma)})(q_\ell^{(\gamma)})^{k_\ell}.
  \label{eq:tgamma}
\end{align}

\subsection{Why the output is diagonal}

Suppose $\Gamma$ is covariant under independent phase rotations in every
mode.  Since $\Phi_0$ is invariant under those rotations, so is
$\Gamma(\Phi_0)$.  The joint eigenspaces of the independent number operators
$(\widehat N_1,\ldots,\widehat N_r)$ are one-dimensional.  Therefore
\begin{equation}\label{eq:P-output}
  \Gamma(\Phi_0)
  =\sum_{\mathbf k}P_{\mathbf k}\ketbra{\mathbf k},
  \qquad
  c:=\sum_{\mathbf k}P_{\mathbf k}\leq1.
\end{equation}
Independent phase covariance is stronger than covariance under only the
total phase rotation; the latter has degenerate eigenspaces and would not by
itself force diagonality in the joint number basis.

\section{The input from Lemma D.4}

Lemma D.4 says that for every bounded product weight
\[
  W(\widehat{\mathbf N})
  =\prod_{\ell=1}^r w_\ell(\widehat N_\ell),
\]
with each $w_\ell$ nonnegative and decreasing,
\begin{equation}\label{eq:D4}
  \sum_{\mathbf k}P_{\mathbf k}w_{\mathbf k}
  \leq
  c\sum_{\mathbf k}t_{\mathbf k}^{(\gamma)}w_{\mathbf k},
  \qquad
  w_{\mathbf k}=\prod_\ell w_\ell(k_\ell).
\end{equation}
This is a conic inequality because it includes the output mass $c$.

Lemma D.4 ultimately depends on the Gu\c{t}\u{a}--Matsumoto least-noise
stochastic-order theorem through Lemma D.3.  That external theorem is not
invoked anew in D.5.  The choice of weight, tangent calculation, and fidelity
argument below are all proved within the manuscript.

\section{Statement of Lemma D.5}

\begin{lemmaD}
Under the assumptions and notation of Lemma D.4,
\begin{equation}\label{eq:D5}
  \B(\Gamma(\Phi_0),\Phi_0)
  \leq\prod_{\ell=1}^r b_\gamma(q_\ell).
\end{equation}
Equality is attained by the product quantum-limited amplifier.
\end{lemmaD}

\section{The tangent inequality}

The scalar function $x\mapsto\sqrt{x}$ is concave, so its tangent line at
$y>0$ lies above its graph.  In the scaled form needed here, for
$x\geq0$ and $s,y>0$,
\begin{equation}\label{eq:tangent-scalar}
  \sqrt{xs}
  \leq
  \sqrt{ys}+\frac12\sqrt{\frac{s}{y}}(x-y).
\end{equation}
There is also a direct identity for the gap:
\begin{equation}\label{eq:tangent-gap}
  \sqrt{ys}+\frac12\sqrt{\frac{s}{y}}(x-y)-\sqrt{xs}
  =\frac12\sqrt{\frac{s}{y}}(\sqrt{x}-\sqrt{y})^2\geq0.
\end{equation}

Apply \eqref{eq:tangent-scalar} separately at every $\mathbf k$ with
\[
  x=P_{\mathbf k},
  \qquad y=t_{\mathbf k}^{(\gamma)},
  \qquad s=t_{\mathbf k}.
\]
All thermal masses are strictly positive.  Define the tangent slope, apart
from the common factor $1/2$, by
\begin{equation}\label{eq:weight}
  w_{\mathbf k}
  :=\sqrt{\frac{t_{\mathbf k}}{t_{\mathbf k}^{(\gamma)}}}.
\end{equation}
Summing the pointwise tangent inequalities gives
\begin{equation}\label{eq:tangent-sum}
  \B(P,t)
  \leq
  \B(t^{(\gamma)},t)
  +\frac12\sum_{\mathbf k}w_{\mathbf k}
      (P_{\mathbf k}-t_{\mathbf k}^{(\gamma)}).
\end{equation}
This is the ``tangent inequality for the square root'' in the paper.  There
is no hidden operator inequality: it is the elementary scalar inequality
\eqref{eq:tangent-scalar}, applied after phase covariance has made all states
diagonal.

The infinite sums are legitimate.  The weights will be bounded, so the two
weighted sums in the correction term are absolutely convergent; the
fidelity sums are finite by Cauchy--Schwarz.

\section{Why the tangent slope is an admissible stochastic weight}

Using \eqref{eq:t}--\eqref{eq:tgamma}, the weight factorizes:
\begin{align}
  w_{\mathbf k}
  &=\prod_{\ell=1}^r w_\ell(k_\ell),
  \label{eq:weight-product}\\
  w_\ell(k)
  &=\sqrt{\frac{1-q_\ell}{1-q_\ell^{(\gamma)}}}
    \left(\frac{q_\ell}{q_\ell^{(\gamma)}}\right)^{k/2}.
  \label{eq:one-weight}
\end{align}
Moreover,
\[
  q_\ell^{(\gamma)}-q_\ell
  =\frac{(\gamma-1)(1-q_\ell)}{\gamma}>0.
\]
Hence $0<q_\ell/q_\ell^{(\gamma)}<1$, and every $w_\ell(k)$ is bounded,
positive, and decreasing.  Thus \eqref{eq:weight-product} belongs exactly to
the product class controlled by Lemma D.4.

This matching is the central design of the proof: the derivative of
fidelity at the quantum-limited output happens to be a decreasing product
observable.

\section{Proof walkthrough}

\subsection*{Step 1: apply Lemma D.4 to the tangent slope}

For the weight \eqref{eq:weight}, equation \eqref{eq:D4} gives
\begin{equation}\label{eq:weighted-order}
  \sum_{\mathbf k}w_{\mathbf k}P_{\mathbf k}
  \leq
  c\sum_{\mathbf k}w_{\mathbf k}t_{\mathbf k}^{(\gamma)}.
\end{equation}
Because $c\leq1$ and every summand is nonnegative,
\begin{equation}\label{eq:correction-negative}
  \sum_{\mathbf k}w_{\mathbf k}
      (P_{\mathbf k}-t_{\mathbf k}^{(\gamma)})
  \leq(c-1)
      \sum_{\mathbf k}w_{\mathbf k}t_{\mathbf k}^{(\gamma)}
  \leq0.
\end{equation}
This also explains why allowing trace-decreasing maps causes no problem:
lost mass only makes the tangent correction smaller.

\subsection*{Step 2: return to the fidelity}

Insert \eqref{eq:correction-negative} into \eqref{eq:tangent-sum}:
\begin{equation}\label{eq:B-reference}
  \B(P,t)\leq\B(t^{(\gamma)},t).
\end{equation}
By \eqref{eq:classical-B}, the left side is
$\B(\Gamma(\Phi_0),\Phi_0)$.  By multiplicativity and
\eqref{eq:bgamma},
\begin{equation}\label{eq:factorize-B}
  \B(t^{(\gamma)},t)
  =\prod_{\ell=1}^r
    \B(\tau_{q_\ell^{(\gamma)}},\tau_{q_\ell})
  =\prod_{\ell=1}^r b_\gamma(q_\ell).
\end{equation}
Equations \eqref{eq:B-reference} and \eqref{eq:factorize-B} prove
\eqref{eq:D5}.

\subsection*{Step 3: check sharpness}

The product quantum-limited amplifier is trace preserving and has the
required displacement and phase covariance.  It maps
\[
  \Phi_0\longmapsto\tau^{(\gamma)},
\]
so its fidelity is exactly the right side of \eqref{eq:factorize-B}.
Therefore the bound is attained.

\section{A one-mode numerical illustration}

Take $q=1/2$ and $\gamma=2$.  Then $q^{(\gamma)}=3/4$, and
\[
  b_2(1/2)
  =\frac{\sqrt{(1/2)(1/4)}}{1-\sqrt{3/8}}
  \approx0.912.
\]
The tangent weight is
\[
  w_k=\sqrt2\left(\sqrt{\frac23}\right)^k,
\]
which decreases with photon number.  The stochastic result says that no
admissible output can have a larger expectation of this low-photon-number
witness than the quantum-limited distribution, after accounting for its
trace.  The tangent inequality then converts precisely that one linear
comparison into the nonlinear fidelity bound.

\section{Imported facts versus the D.5 argument}

\begin{center}
\begin{tabular}{p{0.49\textwidth}p{0.37\textwidth}}
\hline
Fact & Status\\
\hline
Vacuum-idler amplifier is least noisy in photon-number stochastic order &
Imported from Gu\c{t}\u{a}--Matsumoto and used through Lemma D.3\\
One-mode conic inequality and its tensorization & Proved in Lemmas D.3 and
D.4\\
Scalar tangent inequality \eqref{eq:tangent-scalar} & Proved directly; its
nonnegative gap is \eqref{eq:tangent-gap}\\
Factorization and monotonicity of the tangent weight & Checked directly from
the thermal probabilities\\
Sharp fidelity bound and equality case & Proved in Lemma D.5\\
\hline
\end{tabular}
\end{center}

\section{Dependencies and role in the paper}

The full logical chain is
\[
  \text{Lemma D.1}
  +\text{Gu\c{t}\u{a}--Matsumoto}
  \longrightarrow\text{Lemma D.3}
  \longrightarrow\text{Lemma D.4}
  \longrightarrow\text{Lemma D.5}.
\]
For Theorem 3.4, the F\o lner reduction first replaces the flat-prior problem
by an exactly displacement-covariant map evaluated at the origin.  Phase
twirling supplies the phase covariance required here and cannot decrease the
origin fidelity.  Lemma D.5 then gives the upper bound
\[
  \prod_{i<j}b_\gamma(q_{ij})=\mathrm B_{\mathrm{orb}}(\gamma,p).
\]
The product quantum-limited amplifier gives the matching lower bound.

The hybrid converse later uses Lemma D.4 directly rather than D.5: its
branchwise stochastic witness is combined with a separate classical Gaussian
witness.  Thus D.5 is the final step for the purely orbital theorem, while
D.4 is the reusable stochastic component.

The minimal prerequisites are root fidelity for commuting states,
multiplicativity, thermal distributions, phase covariance, and Lemma D.4.
The only calculus idea is the tangent bound for a concave square root;
identity \eqref{eq:tangent-gap} verifies it algebraically.  The specialized
least-noise theorem enters earlier through Lemma D.3, not in the tangent
argument itself.

\end{document}
